%%%%%%%%%%%%%%%%%%%%%%%%%%%%%%%%%%%%%%%%%
% Twenty Seconds Resume/CV
% LaTeX Template
% Version 1.1 (8/1/17)
%
% This template has been downloaded from:
% http://www.LaTeXTemplates.com
%
% Original author:
% Carmine Spagnuolo (cspagnuolo@unisa.it) with major modifications by
% Vel (vel@LaTeXTemplates.com)
%
% License:
% The MIT License (see included LICENSE file)
%
%%%%%%%%%%%%%%%%%%%%%%%%%%%%%%%%%%%%%%%%%

%----------------------------------------------------------------------------------------
%	PACKAGES AND OTHER DOCUMENT CONFIGURATIONS
%----------------------------------------------------------------------------------------

\documentclass[letterpaper]{twentysecondcv} % a4paper for A4

%----------------------------------------------------------------------------------------
%	 PERSONAL INFORMATION
%----------------------------------------------------------------------------------------

% If you don't need one or more of the below, just remove the content leaving the command, e.g. \cvnumberphone{}

\profilepic{profile.jpg} % Profile picture

\cvname{David Abella $\; \; $ Bujalance} % Your name
\cvjobtitle{Physicist. Researcher in Computational Social Science.} % Job title/career

\cvdate{} % Date of birth
\cvaddress{Barcelona, Spain} % Short address/location, use \newline if more than 1 line is required
\cvnumberphone{} % Phone number
\cvsite{\href{https://davidabbu.github.io}{davidabbu.github.io}} % Personal website
\cvmail{davidabbu@gmail.com} % Email address

\hypersetup{hidelinks}

%----------------------------------------------------------------------------------------

\begin{document}

%----------------------------------------------------------------------------------------
%	 ABOUT ME
%----------------------------------------------------------------------------------------

\aboutme{}

%----------------------------------------------------------------------------------------
%	 SKILLS
%----------------------------------------------------------------------------------------

% Skill bar section, each skill must have a value between 0 an 6 (float)
\skills{{Python /5.5},{Julia /5.5},{Fortran /5.5},{Office/5.9},{Latex /5.0}}

%----------------------------------------------------------------------------------------
%	 PAGE 1
%----------------------------------------------------------------------------------------

\makeprofile % Print the sidebar

\section{Interests}

\medskip

As a physicist, my expertise lies in statistical mechanics, network science and data analysis, a powerful toolkit to decipher complex problems. I apply these skills to social and urban challenges, such as housing markets, urban segregation, collective behavior and the governance of online platforms like Wikipedia.

\bigskip

\section{Work Experience}

\begin{twenty} % Environment for a list with descriptions
	\twentyitem{2025-Now}{Researcher in Computational Social Science at Eurecat}{}{Eurecat, Centre Tecnològic de Catalunya. Projects: DEM-Debate (Wikipedia editing and content moderation during the 2024 European Parliament election, with Wikimedia Europe and IViR, University of Amsterdam); AI4Bus (AI forecasting of passenger flows for Barcelona's metropolitan bus network, with AMB); deployEMDS (deployment of the common European Mobility Data Space). \smallskip}
	\twentyitem{2023-2024}{Research Technician at IFISC (CSIC)}{}{Research associated with the projects Nou Lloguer and Next4Mob to continue renting market data analysis and modelling. \smallskip}
	\twentyitem{2022-2023}{Young researchers scholarship by SOIB (UIB)}{}{Research focused on different agent-based models of human dynamics. \smallskip}
	\twentyitem{2022}{Research Technician at IFISC (CSIC)}{}{Research associated with the projects APASOS-CSIC and Nou Lloguer for the renting market data analysis. \smallskip}
	\twentyitem{2021}{Predoctoral contract at IFISC}{}{8-month contract to continue the study of aging effects in threshold models with application in sociotechnical systems. Focused on the study of aging in the Noisy Schelling segregation model. \smallskip}
\end{twenty}

\medskip

\section{Education}

\medskip

\begin{twenty} % Environment for a list with descriptions
    \twentyitem{2021-2024}{PhD {\normalfont Physics of Complex Systems} at the Institute for Cross-Disciplinary Physics and Complex Systems (IFISC, UIB-CSIC)}{}{Thesis: \textit{Aging and memory effects in social and economic dynamics}. Supervisors: Maxi San Miguel and José J. Ramasco. Awarded the doctoral thesis prize of the UIB Chair of Urban Studies (2024).}

    \twentyitem{2019-2020}{M.Sc {\normalfont Physics of Complex Systems} at Universitat de les Illes Balears}{}{Complex systems and computational social science, applying physics to interdisciplinary problems.
	\smallskip
	\begin{itemize}
	    \item Master thesis: Perturbations propagation in self-organized patterns.
	\end{itemize}
	}

    \twentyitem{2015-2019}{B.Sc {\normalfont Physics} at Universitat de Barcelona}{}{Minor in Fundamental Physics.
	\smallskip
	\begin{itemize}
	    \item Bachelor thesis: Many-body contributions in water clusters (published in ACS Nano).
	\end{itemize}
	}
\end{twenty}

%----------------------------------------------------------------------------------------
%	 PAGE 2
%----------------------------------------------------------------------------------------

\newpage
\setcounter{colorCounter}{0} % Keep section titles green on every page

\makeprofile % Print the sidebar

\section{Scientific publications}

\begin{twenty}
    \twentyitem{2026}{Computational analysis of Wikipedia editorial activity during the 2024 European Parliament election}{}{D. Abella Bujalance, D. Laniado. Software and data, \href{https://zenodo.org/records/22111198}{Zenodo}.}

    \twentyitem{2025}{Exploring the spatial segmentation of housing markets from online listings}{}{EPJ Data Science 14, 34.}

    \twentyitem{2024}{Ordering dynamics and aging in the symmetrical threshold model}{}{New Journal of Physics 26, 013033.}

    \twentyitem{2023}{Aging in binary-state models: the Threshold model for complex contagion}{}{Physical Review E 107, 024101.}

    \twentyitem{2023}{Many-body contributions in water nanoclusters}{}{ACS Nano 17 (3), 1959-1964.}

    \twentyitem{2023}{Unraveling higher-order dynamics in collaboration networks}{}{Preprint, \href{https://doi.org/10.48550/arXiv.2306.17521}{arXiv:2306.17521}.}

    \twentyitem{2022}{Aging effects in Schelling segregation model}{}{Scientific Reports 12, 19376.}
\end{twenty}

\section{Awards}

\begin{twenty}
    \twentyitem{Nov-2024}{Doctoral thesis prize, UIB Chair of Urban Studies (2nd edition)}{}{Prize of 1,500\,€ for the thesis \textit{Aging and memory effects in social and economic dynamics}.}
\end{twenty}

\section{Collaboration and outreach events}

\begin{twenty} % Environment for a list with descriptions
	\twentyitem{2024}{Pint of Science 2024}{}{Talk for a general audience on the physics of social and urban dynamics. \smallskip}
	\twentyitem{Sep-2023}{Associate researcher at Queen Mary University of London (QMUL)}{}{iMOVE research stay with Vincenzo Nicosia to study modeling techniques for spatial segregation. \smallskip}
	\twentyitem{2022-2023}{Poster Party}{}{Member of the Poster Party committee organizing the outreach event at IFISC. \smallskip}
	\twentyitem{Oct-2022}{Conference on Complex Systems 2022}{}{Volunteer and collaborator at the Conference on Complex Systems CCS 2022 in Palma. \smallskip}
\end{twenty}

%----------------------------------------------------------------------------------------
%	 PAGE 3
%----------------------------------------------------------------------------------------

\newpage
\setcounter{colorCounter}{0} % Keep section titles green on every page

\makeprofile % Print the sidebar

\section{Scientific presentations}

\begin{twenty}
    \twentyitem{Nov-2026}{Oral presentation "AI4Bus: Artificial intelligence applied to the metropolitan bus network" at the XII Symposium on Urban Mobility Challenges, CARNET (Barcelona - Spain)}{}{}

    \twentyitem{Feb-2026}{Presentation of the DEM-Debate results at the project final event, European Parliament (Brussels - Belgium)}{}{}

    \twentyitem{Sep-2025}{Oral presentation "Exploring Wikipedia community practices during the 2024 European Parliament election" at CCS2025 (Siena - Italy)}{}{}

    \twentyitem{May-2025}{Presentation "Exploring Wikipedia community practices during the 2024 European Parliament election" at Wiki Workshop 2025 (online)}{}{}

    \twentyitem{Feb-2024}{Oral presentation "Assessing the real estate spatial segmentation" at CS3-24 (Barcelona - Spain)}{}{}

    \twentyitem{July-2023}{Oral presentation "Aging in binary state dynamics: Ordering in threshold models" at NetSci23 (Vienna - Austria)}{}{}

    \twentyitem{May-2023}{Oral presentation "Aging in binary state dynamics: Ordering in threshold models" at CS3-23 (Santander - Spain)}{}{}

    \twentyitem{Oct-2022}{Oral presentation "Aging in binary-state models: Threshold model of complex contagion" at CCS2022 (Palma - Spain)}{}{}

    \twentyitem{Sep-2022}{Poster presentation "Aging in binary-state models" at II Summer school on Complex Socio-Technical systems (Palma - Spain)}{}{}

    \twentyitem{Aug-2022}{Poster presentation "Aging in binary-state models" at X Summer school on Statistical physics of Complex Systems (Palma - Spain)}{}{}

    \twentyitem{May-2022}{Poster presentation "Aging in Complex contagion" at FISES (Zaragoza - Spain)}{}{}

    \twentyitem{Mar-2022}{Oral presentation "Aging effects in Schelling segregation model" at COMSOTEC Workshop (Tarragona - Spain)}{}{}

    \twentyitem{Oct-2021}{Oral presentation "Aging effects in Schelling segregation model" at CCS2021 (Lyon - France)}{}{}
\end{twenty}

\section{Supplementary Education}

\begin{twenty}
    \twentyitem{Oct-2026}{Agile training by Kairós (at Eurecat)}{}{}
    \twentyitem{2023}{Participation at the Complexity72h Workshop}{}{}
    \twentyitem{2022}{Summer school about Artificial Intelligence "Retos de la IA" by AIHUB-CSIC}{}{}
    \twentyitem{2022}{X Summer school on Statistical physics of Complex Systems by GEFENOL (at IFISC)}{}{}
    \twentyitem{2022}{II Summer school on Complex Socio-Technical systems by COMSOTEC (at IFISC)}{}{}
    \twentyitem{2021}{Introduction course to Artificial Intelligence at UBICS}{}{}
\end{twenty}

\end{document}
